\subsubsection{Transaction Cost Analysis (TCA)}

\begin{definition} \hlt{TCA}\\
Transaction cost analysis compares realised execution prices with relevant benchmarks and explains the gap by market context, order design, broker behaviour, and trading style.
\end{definition}

\begin{principle} \hlt{Pre-, Live-, and Post-Trade Consistency}\\
The same indicator should have a pre-trade expectation, a live monitored value, and a post-trade assessment. A bid-ask-spread indicator, for example, can move from expected average, to updated live likelihood, to realised influence on performance.
\end{principle}

\begin{taxonomy} \hlt{Performance Attribution Segmentation}
\begin{enumerate}[label=\roman*.]
\setlength{\itemsep}{0pt}
\item \hlt{immediate price}: the reference available when the order decision is made;
\item \hlt{planned price}: the price expected from the risk-control trajectory;
\item \hlt{unexpected market move}: exogenous price variation during execution;
\item \hlt{trading efficiency}: quality of liquidity-seeking tactics relative to the plan;
\item \hlt{market impact}: price movement caused by the order itself.
\end{enumerate}
This segmentation is useful only when the benchmark, horizon, and side are explicit.
\end{taxonomy}

\begin{caveat} \hlt{TCA Uses Wider Samples Than Post-Trade Review}\\
Post-trade review explains one execution in detail. TCA usually aggregates months of comparable trades to identify contexts where an algorithm, broker, or parameter set is systematically strong or weak.
\end{caveat}

\begin{taxonomy} \hlt{TCA Context Segmentation}
\begin{enumerate}[label=\roman*.]
\setlength{\itemsep}{0pt}
\item \hlt{liquidity}: spread, queue size, turnover, trade size, free float, and trading rate;
\item \hlt{risk}: volatility, jumps, intraday trend, and index or sector co-movement;
\item \hlt{order design}: size, duration, participation rate, benchmark, and limit constraints;
\item \hlt{market state}: news, auctions, rebalances, corporate events, and venue outages.
\end{enumerate}
\end{taxonomy}

\begin{practitionernote} \hlt{Broker and Algo Selection}\\
The practical output of TCA is a routing table: for a given market context and trading objective, send the next order to the broker or algorithm whose historical performance is best after adjusting for risk and liquidity.
\end{practitionernote}

\begin{principle} \hlt{Best Execution as Auditability}\\
Best execution is not the cheapest commission or the fastest fill in isolation. It is the best auditable outcome given the client's objective, limit, urgency, available venues, and information risk.
\end{principle}

\begin{taxonomy} \hlt{Broker-Conflict Segmentation}
\begin{enumerate}[label=\roman*.]
\setlength{\itemsep}{0pt}
\item \hlt{routing inducements}: rebates, payments for order flow, reciprocal flow, and affiliate venues;
\item \hlt{information misuse}: front-running, inappropriate exposure, and prearranged trades;
\item \hlt{allocation}: assigning favourable fills to preferred accounts;
\item \hlt{activity incentives}: commissions, churning, soft dollars, and commission recapture;
\item \hlt{principal conflict}: internalisation, dual trading, and dealer inventory.
\end{enumerate}
TCA and audit trails convert these conflicts from stories into measurable routing and fill-quality checks.
\end{taxonomy}

\begin{process} \hlt{Block Motive Audit}
\begin{enumerate}[label=\roman*.]
\setlength{\itemsep}{0pt}
\item identify economic reason for trade: funding, diversification, hedging, benchmark change, or speculation;
\item test whether the requested size and timing are consistent with that reason;
\item compare the initiator's alternatives, constraints, and past behaviour;
\item price the trade only after separating utilitarian urgency from possible information.
\end{enumerate}
The audit is not moral judgement; it is adverse-selection control.
\end{process}

\begin{caveat} \hlt{Block Trade Reporting Trade-Off}\\
Fast public reporting improves transparency and protects the regular market, but it can hurt dealers that still need to unwind facilitated inventory. Delayed reporting protects facilitation, but gives informed block traders more room to exploit stale prices.
\end{caveat}
